Việc huấn luyện mạng nơ-ron phụ thuộc vào cách tín hiệu và đạo hàm biến đổi qua các lớp. Ba vấn đề cần phân biệt là: độ nhạy cục bộ của hàm kích hoạt; độ nhạy của hợp các ánh xạ trong toàn mạng; và thang đo thống kê của tín hiệu tại thời điểm khởi tạo. Phân tích sau xây dựng lần lượt ba nội dung này, rồi suy ra các quy tắc khởi tạo theo số chiều đầu vào và đầu ra. Các quy tắc được trình bày
với phạm vi giả thiết cụ thể, thay vì xem chúng là bảo đảm hội tụ cho mọi mạng và mọi bài toán.

\section{Vùng bão hòa của hàm kích hoạt}
\label{kt:sec:bao-hoa}

\subsection{Tiền kích hoạt và khái niệm bão hòa}
Xét mạng nơ-ron truyền thẳng có $L$ lớp (không tính lớp đầu vào $l = 0$):
\begin{align}
    \mathbf a^{[0]}
    &=\mathbf x,\notag\\
    \mathbf z^{[l]}
    &=\mathbf W^{[l]}\mathbf a^{[l-1]}+
    \mathbf b^{[l]},\label{kt:eq:forward-z}\\
    \mathbf a^{[l]}&=\sigma^l(\mathbf z^{[l]}),
    \qquad l=1,\ldots,L.
    \label{kt:eq:forward-a}
\end{align}
Trong đó, các đại lượng toán học được định nghĩa chi tiết như sau:
\begin{itemize}
    \item $\mathbf{x}, \mathbf a^{[0]} \in \mathbb{R}^{n_0}$: Véc-tơ tín hiệu đầu vào của mạng nơ-ron.
    
    \item $n_l \in \mathbb{N}$: Số lượng nơ-ron tại lớp thứ $l$.
    
    \item $\mathbf{W}^{[l]} \in \mathbb{R}^{n_l \times n_{l-1}}$: Ma trận trọng số liên kết các tín hiệu từ lớp $l-1$ truyền đến lớp $l$.
    
    \item $\mathbf{b}^{[l]} \in \mathbb{R}^{n_l}$: Véc-tơ hệ số chệch của lớp thứ $l$.
    
    \item $\mathbf{z}^{[l]} \in \mathbb{R}^{n_l}$: Véc-tơ tiền kích tại lớp thứ $l$.
    
    \item $\mathbf{a}^{[l]} \in \mathbb{R}^{n_l}$: Véc-tơ tín hiệu đầu ra của lớp thứ $l$.
    
    \item $\sigma^l$: Hàm kích hoạt tại lớp $l$. Phép toán này tác động độc lập lên từng phần tử của véc-tơ.
\end{itemize}
Cụ thể,
\begin{equation}
 z_i^{[l]}=\sum_{j=1}^{n_{l-1}}W_{ij}^{[l]}a_j^{[l-1]}+b_i^{[l]}.
 \label{kt:eq:neuron}
\end{equation}
Chỉ số $i$ xác định nơ-ron nhận tín hiệu, còn $j$ xác định nơ-ron
ở lớp trước. Mỗi nơ-ron nhận một độ lệch $b_i^{[l]}$ ngoài tổng.
Cả trọng số và độ lệch thường được tối ưu. Với các lớp ẩn, đầu vào
$\mathbf a^{[l-1]}$ cũng thay đổi gián tiếp khi tham số những lớp
trước được cập nhật; chỉ dữ liệu gốc $\mathbf x$ được giữ cố định
trong một lần đánh giá mạng.

Đối với hàm trơn, khai triển tại $z$ cho
\begin{equation}
 \phi(z+h)-\phi(z)=\phi'(z)h+o(|h|),\qquad h\to0.
 \label{kt:eq:local-sensitivity}
\end{equation}
Vì vậy $|\phi'(z)|$ đo độ nhạy bậc nhất của đầu ra đối với nhiễu
nhỏ ở tiền kích hoạt. Với các hàm bị chặn như logistic và tang
hyperbol, đạo hàm giảm về không khi đi xa về hai phía; đây là
cơ sở của khái niệm bão hòa. Không có một ngưỡng hữu hạn duy nhất
phân chia tuyệt đối vùng bão hòa và không bão hòa.

\begin{definition}[Bão hòa theo ngưỡng đạo hàm tương đối]
Với một trong hai hàm logistic hoặc tang hyperbol, đặt
$M_\phi=\sup_{z\in\mathbb R}|\phi'(z)|>0$. Cho
$\varepsilon\in(0,1)$, định nghĩa
\begin{equation}
 d_\phi(z)=\frac{|\phi'(z)|}{M_\phi},\qquad
 \mathcal S_{\phi,\varepsilon}
   =\{z\in\mathbb R:d_\phi(z)\leq\varepsilon\}.
 \label{kt:eq:saturation-set}
\end{equation}
\end{definition}
Trong đó, $d_\phi$ là độ nhạy tương đối, còn
$\mathcal S_{\phi,\varepsilon}$ là vùng được quy ước là bão hòa ở
ngưỡng $\varepsilon$. Định nghĩa này phục vụ so sánh định lượng;
nó không khẳng định mạng ngừng học ngay khi một nơ-ron vượt ngưỡng.

\subsection{Hàm logistic: đạo hàm và ngưỡng bão hòa}

Đặt
\begin{equation}
 \sigma(z)=\frac{1}{1+\exp(-z)}.
 \label{kt:eq:sigmoid}
\end{equation}
Lấy đạo hàm của nghịch đảo và dùng
$1-\sigma(z)=\exp(-z)/(1+\exp(-z))$, ta được
\begin{align}
 \sigma'(z)
 &=\frac{\exp(-z)}{(1+\exp(-z))^2}
   =\sigma(z)(1-\sigma(z))\notag\\
 &=\frac14-\left(\sigma(z)-\frac12\right)^2
   =\frac{1}{4\cosh^2(z/2)}.
 \label{kt:eq:sigmoid-prime}
\end{align}
Trong đó, $\cosh t=(\exp(t)+\exp(-t))/2$. Dạng bình phương cho
$0<\sigma'(z)\leq1/4$, với dấu bằng ở $z=0$; dạng chứa $\cosh$
cho thấy đạo hàm là hàm chẵn và giảm theo $|z|$. Do đó
$M_\sigma=1/4$ và
\begin{align}
 d_\sigma(z)\leq\varepsilon
 &\Longleftrightarrow\frac{1}{\cosh^2(z/2)}\leq\varepsilon\notag\\
 &\Longleftrightarrow\cosh(|z|/2)\geq\varepsilon^{-1/2}\notag\\
 &\Longleftrightarrow |z|\geq
 c_\sigma(\varepsilon):=2\operatorname{arcosh}(\varepsilon^{-1/2}).
 \label{kt:eq:sigmoid-threshold}
\end{align}
Hàm $\operatorname{arcosh}(v)=\log(v+\sqrt{v^2-1})$, $v\geq1$,
là hàm ngược của $\cosh$ trên nửa trục không âm. Các phép tương
đương sử dụng tính tăng của $\cosh$ trên miền này.

Chẳng hạn, $\sigma'(2)\approx0{,}104994$ và
$\sigma'(4)\approx0{,}017663$. Đây là các giá trị làm tròn,
không nên thay chúng thành cận dưới chính xác $0{,}105$ và
$0{,}018$. Ngay cả tại điểm có đạo hàm lớn nhất, mỗi phép nhân
riêng với $\sigma'$ vẫn có hệ số không quá $1/4$; tính ổn định của
mạng còn phụ thuộc các ma trận trọng số.

\subsection{Hàm tang hyperbol và sự khác biệt về thang đo}

Hàm tang hyperbol có biểu thức
\begin{equation}
 \tanh z=\frac{\exp(z)-\exp(-z)}{\exp(z)+\exp(-z)}
        =2\sigma(2z)-1.
 \label{kt:eq:tanh}
\end{equation}
Quy tắc dây chuyền hoặc đạo hàm của thương cho
\begin{equation}
 (\tanh z)'=4\sigma'(2z)=1-\tanh^2z
           =\frac1{\cosh^2z}.
 \label{kt:eq:tanh-prime}
\end{equation}
Vì $\cosh z\geq1$, đạo hàm thuộc $(0,1]$ và đạt cực đại bằng một
tại $z=0$. Như vậy $M_{\tanh}=1$ và
\begin{equation}
 \mathcal S_{\tanh,\varepsilon}
 =\{z:|z|\geq c_{\tanh}(\varepsilon)\},\qquad
 c_{\tanh}(\varepsilon)=\operatorname{arcosh}(\varepsilon^{-1/2}).
 \label{kt:eq:tanh-threshold}
\end{equation}
Trong đó, $c_{\tanh}$ được suy ra bằng các bước như
\eqref{kt:eq:sigmoid-threshold}. Với cùng ngưỡng tương đối,
$c_\sigma=2c_{\tanh}$; vì vậy việc so sánh hai hàm phải chỉ rõ
đang xét cùng đầu vào, cùng đạo hàm tuyệt đối hay cùng mức suy giảm
tương đối.

\begin{table}[htbp]
\centering
\caption{Ngưỡng bão hòa tương đối, làm tròn đến sáu chữ số thập phân.}
\label{kt:tab:threshold}
\begin{tabular}{@{}ccc@{}}
\toprule
$\varepsilon$&$c_\sigma(\varepsilon)$&$c_{\tanh}(\varepsilon)$\\
\midrule
$0{,}10$&$3{,}636893$&$1{,}818446$\\
$0{,}05$&$4{,}356544$&$2{,}178272$\\
$0{,}01$&$5{,}986446$&$2{,}993223$\\
\bottomrule
\end{tabular}
\end{table}

Các giá trị cần chú ý là
$(\tanh)'(1)\approx0{,}419974$,
$(\tanh)'(2)\approx0{,}070651$ và
$(\tanh)'(2{,}5)\approx0{,}026592$.
Đạo hàm cực đại của $\tanh$ lớn gấp bốn lần logistic, nhưng điều
này chưa đủ chứng minh mọi mạng dùng $\tanh$ có građien ổn định hơn:
đối số, trọng số và hàm mất mát có thể khác nhau.

% \begin{figure}[htbp]
%  \centering
%  \includegraphics[width=\linewidth]{hinh/ham_kich_hoat.png}
%  \caption{Hàm kích hoạt và độ nhạy tương đối. Đường ngang ở hình
%  bên phải biểu diễn ngưỡng $\varepsilon=0{,}1$; giao điểm xác định
%  các ngưỡng ở Bảng~\ref{kt:tab:threshold}. Hình được dựng trực tiếp
%  từ công thức, không phải kết quả huấn luyện.}
%  \label{kt:fig:activation}
% \end{figure}

\subsection{Đo lường bão hòa trên một tập dữ liệu}

Với $N$ mẫu đánh giá và $n_l$ nơ-ron ở lớp $l$, đặt
\begin{align}
 \widehat S_{\varepsilon}^{[l]}
 &=\frac1{Nn_l}\sum_{p=1}^N\sum_{i=1}^{n_l}
   \mathbf1\!\left\{d_{\phi_l}(z_{i,p}^{[l]})\leq\varepsilon\right\},
      \label{kt:eq:saturation-fraction}\\
 \widehat D^{[l]}
 &=\frac1{Nn_l}\sum_{p=1}^N\sum_{i=1}^{n_l}
                         d_{\phi_l}(z_{i,p}^{[l]}),
      \label{kt:eq:mean-derivative}\\
 \widehat\rho^{[l]}
 &=\frac1{Nn_l}\sum_{p=1}^N\sum_{i=1}^{n_l}|z_{i,p}^{[l]}|.
      \label{kt:eq:mean-preactivation}
\end{align}
Trong đó, $p$ chỉ mẫu dữ liệu; $\mathbf1\{A\}$ bằng một nếu mệnh
đề $A$ đúng, bằng không nếu sai. Đại lượng $\widehat S$ là tỉ lệ
bão hòa thực nghiệm; $\widehat D$ là độ nhạy tương đối trung bình;
$\widehat\rho$ chỉ là trung bình độ lớn tiền kích hoạt. Chúng phụ
thuộc tập dữ liệu đánh giá và bộ tham số hiện tại.

Với logistic hoặc $\tanh$, viết $c=c_{\phi_l}(\varepsilon)>0$.
Bất đẳng thức $\mathbf1\{|z|\geq c\}\leq|z|/c$ cho
\begin{equation}
 \widehat S_{\varepsilon}^{[l]}
 \leq\min\{1,\widehat\rho^{[l]}/c\}.
 \label{kt:eq:rho-bound}
\end{equation}
Như vậy $\widehat\rho$ nhỏ có thể cho một cận trên hữu ích, nhưng
$\widehat\rho$ lớn chưa buộc tỉ lệ bão hòa lớn. Ví dụ, với ngưỡng
$c=3$, một trăm tiền kích hoạt đều bằng $4$ có
$\widehat\rho=4$, $\widehat S=1$; một tiền kích hoạt bằng $400$
và chín mươi chín giá trị bằng không vẫn có $\widehat\rho=4$,
nhưng $\widehat S=0{,}01$. Vì thế không dùng riêng trung bình độ
lớn để kết luận phần lớn nơ-ron bị bão hòa.

Građien nhỏ cũng có thể do mất mát đã nhỏ, các tín hiệu lỗi triệt
tiêu nhau khi lấy trung bình, hoặc ma trận trọng số làm co một
hướng. Chẩn đoán bão hòa cần quan sát trực tiếp đạo hàm kích hoạt
và phân phối tiền kích hoạt, bên cạnh độ lớn cập nhật tham số.

\subsection{Bão hòa và sai số số học}

Đạo hàm của logistic và $\tanh$ dương tại mọi đối số hữu hạn trong
số học thực. Trong số học dấu phẩy động, đầu ra có thể bị làm tròn
thành đúng $0$, $1$ hoặc $-1$; lúc ấy các biểu thức như
$a(1-a)$ và $1-a^2$ có thể cho không. Bão hòa về mặt giải tích và
giá trị bằng không do làm tròn là hai hiện tượng cần phân biệt.

Một cách đánh giá logistic tránh lấy hàm mũ của số dương lớn là
\begin{equation}
 \sigma(z)=
 \begin{cases}
  (1+\exp(-z))^{-1},&z\geq0,\\
  \exp(z)/(1+\exp(z)),&z<0.
 \end{cases}
 \label{kt:eq:stable-sigmoid}
\end{equation}
Tương tự, đặt $t=\exp(-|z|)$ và $r=\exp(-2|z|)$, ta có
\begin{equation}
 \sigma'(z)=\frac{t}{(1+t)^2},\qquad
 (\tanh z)'=\frac{4r}{(1+r)^2}.
 \label{kt:eq:stable-derivatives}
\end{equation}
Các công thức suy ra từ tính chẵn của đạo hàm và phép nhân tử số,
mẫu số với hàm mũ thích hợp. Chúng tránh một số phép trừ gần nhau
và nguy cơ tràn của hàm mũ, nhưng vẫn có thể xảy ra tràn dưới
khi đạo hàm quá nhỏ. Cải thiện cách tính một đạo hàm không
loại bỏ sự suy giảm thực sự của tích đạo hàm qua nhiều lớp.

\section{Triệt tiêu và bùng nổ građien}
\label{kt:sec:gradient}

\subsection{Các hệ thức lan truyền ngược chính xác}

Giả sử các đạo hàm xuất hiện dưới đây tồn tại tại điểm đánh giá.
Với ReLU, các công thức đạo hàm cổ điển áp dụng ngoài điểm gãy;
giá trị được cài đặt tại điểm gãy là một quy ước riêng.
Với một mẫu dữ liệu, gọi $\mathcal L$ là hàm mất mát vô hướng và đặt
\begin{equation}
 \boldsymbol\delta^{[l]}=\nabla_{\mathbf z^{[l]}}\mathcal L,
 \qquad
 \mathbf D^{[l]}=\operatorname{diag}
                     (\phi_l'(\mathbf z^{[l]})).
 \label{kt:eq:delta-definition}
\end{equation}
Trong đó, $\boldsymbol\delta^{[l]}\in\mathbb R^{n_l}$ là độ nhạy
của mất mát theo tiền kích hoạt; $\mathbf D^{[l]}$ là ma trận đường
chéo chứa các đạo hàm của hàm kích hoạt tại lớp $l$.
Quy tắc dây chuyền cho
\begin{align}
 \boldsymbol\delta^{[L]}
 &=\mathbf D^{[L]}\nabla_{\mathbf a^{[L]}}\mathcal L,\notag\\
 \boldsymbol\delta^{[l]}
 &=\mathbf D^{[l]}(\mathbf W^{[l+1]})^T
                         \boldsymbol\delta^{[l+1]},
                         \qquad 1\leq l<L,\label{kt:eq:delta-recursion}\\
 \nabla_{\mathbf W^{[l]}}\mathcal L
 &=\boldsymbol\delta^{[l]}(\mathbf a^{[l-1]})^T,
 \qquad \nabla_{\mathbf b^{[l]}}\mathcal L
 =\boldsymbol\delta^{[l]}.
 \label{kt:eq:parameter-gradient}
\end{align}
Để kiểm tra công thức theo từng thành phần, dùng
$\partial z_k^{[l+1]}/\partial a_i^{[l]}=W_{ki}^{[l+1]}$ và
$\partial a_i^{[l]}/\partial z_i^{[l]}=\phi_l'(z_i^{[l]})$:
\[
 \delta_i^{[l]}=\phi_l'(z_i^{[l]})
                   \sum_{k=1}^{n_{l+1}}W_{ki}^{[l+1]}\delta_k^{[l+1]}.
\]
Mặt khác, $\partial z_i^{[l]}/\partial W_{ij}^{[l]}=a_j^{[l-1]}$
và $\partial z_i^{[l]}/\partial b_i^{[l]}=1$, nên nhận được
\eqref{kt:eq:parameter-gradient}. Công thức này cho thấy độ lệch
cũng được cập nhật bằng đạo hàm.

Nếu $\mathcal L_B=B^{-1}\sum_{p=1}^B\mathcal L_p$ là mất mát của
một lô $B$ mẫu được đánh giá độc lập qua mạng, thì
\begin{equation}
 \nabla_{\mathbf W^{[l]}}\mathcal L_B
 =\frac1B\sum_{p=1}^B
   \boldsymbol\delta_p^{[l]}(\mathbf a_p^{[l-1]})^T.
 \label{kt:eq:batch-gradient}
\end{equation}
Vì chuẩn Frobenius của tích ngoài bằng tích hai chuẩn Euclid,
\begin{equation}
 \|\nabla_{\mathbf W^{[l]}}\mathcal L_B\|_F
 \leq\frac1B\sum_{p=1}^B
             \|\boldsymbol\delta_p^{[l]}\|_2\,
             \|\mathbf a_p^{[l-1]}\|_2.
 \label{kt:eq:batch-bound}
\end{equation}
Trong đó, $\|\mathbf A\|_F^2=\sum_{i,j}A_{ij}^2$. Cận trên
không phải đẳng thức nói chung, vì các građien của từng mẫu có thể
triệt tiêu nhau. Do đó cần phân biệt građien theo tiền kích hoạt,
građien theo tham số và građien sau khi lấy trung bình theo lô.

\subsection{Tích ma trận Jacobi và điều kiện co giãn}

Đặt
\begin{equation}
 \mathbf J^{[l]}=
 \frac{\partial\mathbf a^{[l]}}{\partial\mathbf a^{[l-1]}}
 =\mathbf D^{[l]}\mathbf W^{[l]},\qquad
 \mathbf g^{[l]}=\nabla_{\mathbf a^{[l]}}\mathcal L.
 \label{kt:eq:layer-jacobian}
\end{equation}
Trong đó, $\mathbf J^{[l]}$ có kích thước $n_l\times n_{l-1}$.
Lặp quy tắc dây chuyền cho
\begin{equation}
 \mathbf g^{[l]}
 =(\mathbf J^{[l+1]})^T\cdots(\mathbf J^{[L]})^T\mathbf g^{[L]}.
 \label{kt:eq:jacobian-product}
\end{equation}
Chuẩn phổ $\|\mathbf A\|_2=\sup_{\|\mathbf v\|_2=1}
\|\mathbf A\mathbf v\|_2$ thỏa tính dưới nhân. Vì vậy
\begin{equation}
 \|\mathbf g^{[l]}\|_2
 \leq\left(\prod_{k=l+1}^{L}
   d_k\|\mathbf W^{[k]}\|_2\right)\|\mathbf g^{[L]}\|_2,
 \qquad d_k=\max_i|\phi_k'(z_i^{[k]})|.
 \label{kt:eq:gradient-product-bound}
\end{equation}
Nếu mọi thừa số $d_k\|\mathbf W^{[k]}\|_2\leq q<1$, thì
\begin{equation}
 \|\mathbf g^{[l]}\|_2\leq q^{L-l}\|\mathbf g^{[L]}\|_2.
 \label{kt:eq:gradient-contraction}
\end{equation}
Đây là điều kiện đủ cho sự suy giảm tương đối theo độ sâu. Muốn
kết luận građien tiến về không trong một họ mạng sâu dần, còn phải
kiểm soát građien tại đầu ra của họ mạng ấy.

Ngược lại, một tích cận trên lớn chưa chứng minh građien thực tế
tăng. Ví dụ
\begin{equation}
 \mathbf A=\begin{pmatrix}2&0\\0&1/2\end{pmatrix},\qquad
 \mathbf B=\begin{pmatrix}1/2&0\\0&2\end{pmatrix}
 \quad\Longrightarrow\quad
 \|\mathbf A\|_2\|\mathbf B\|_2=4,
 \quad\mathbf A\mathbf B=\mathbf I.
 \label{kt:eq:matrix-counterexample}
\end{equation}
Các hướng được khuếch đại ở hai bước khác nhau. Với ma trận vuông,
một điều kiện đủ mạnh hơn cho khuếch đại mọi hướng là
$s_{\min}(\mathbf J^{[k]})\geq c>1$, khi đó
\begin{equation}
 \|\mathbf g^{[l]}\|_2\geq c^{L-l}\|\mathbf g^{[L]}\|_2.
 \label{kt:eq:gradient-expansion}
\end{equation}
Ở đây $s_{\min}$ là giá trị kỳ dị nhỏ nhất, và cận dưới suy ra bằng
áp dụng liên tiếp $\|\mathbf J^T\mathbf v\|_2\geq
s_{\min}(\mathbf J)\|\mathbf v\|_2$. Giả thiết này chỉ là điều
kiện đủ; hiện tượng bùng nổ cũng có thể xảy ra trên một số hướng.

Đối với mạng logistic, $d_k\leq1/4$, nhưng không được thay
$\|\mathbf W^{[k]}\|_2$ bằng độ lớn của một trọng số điển hình.
Ma trận có nhiều phần tử nhỏ vẫn có thể có chuẩn phổ lớn. Cách
phân tích qua các giá trị kỳ dị là một nội dung được nhấn mạnh
trong nghiên cứu của Glorot và Bengio~\cite{glorot2010}.

\subsection{Ví dụ vô hướng: suy giảm và khuếch đại tại cùng mức kích hoạt}

Xét $L$ lớp, mỗi lớp một nơ-ron logistic, đầu vào $x=1$ và mất
mát $\mathcal L=\tfrac12(a^{[L]}-y)^2$ với $y=0$. Chọn
\begin{equation}
 w^{[l]}=w>0,\qquad b^{[1]}=-w,
 \qquad b^{[l]}=-w/2\quad(l=2,\ldots,L).
 \label{kt:eq:scalar-construction}
\end{equation}
Các $w^{[l]}$ vẫn là những tham số riêng; $w$ chỉ là giá trị chung
được gán cho chúng tại điểm đánh giá. Tại mẫu đang xét,
$z^{[1]}=w-w=0$, nên $a^{[1]}=1/2$.
Nếu $a^{[l-1]}=1/2$ thì $z^{[l]}=w/2-w/2=0$.
Bằng quy nạp, mọi $a^{[l]}=1/2$ và mọi đạo hàm kích hoạt bằng
$1/4$. Các độ lệch này là các giá trị cố định tại điểm đánh giá;
không lấy đạo hàm của chúng theo trọng số khi tính đạo hàm riêng.

Từ \eqref{kt:eq:delta-recursion},
\begin{align}
 \delta^{[L]}&=(a^{[L]}-y)\sigma'(z^{[L]})=\frac18,\notag\\
 \delta^{[l]}&=\frac{w}{4}\delta^{[l+1]},\notag\\
 \delta^{[1]}&=\frac18\left(\frac{w}{4}\right)^{L-1},\qquad
 \frac{\partial\mathcal L}{\partial w^{[1]}}=\delta^{[1]},\qquad
 \frac{\partial\mathcal L}{\partial w^{[L]}}=\frac1{16}\quad(L\geq2).
 \label{kt:eq:scalar-gradient}
\end{align}
Thừa số $w/4$ là mức khuếch đại qua một chuyển tiếp trong lan
truyền ngược. Với $L=4$ (ba lớp ẩn và một lớp đầu ra), ta có
\begin{table}[htbp]
\centering
\caption{Đạo hàm chính xác tại mẫu $x=1$, $y=0$ của ví dụ vô hướng.}
\label{kt:tab:scalar}
\begin{tabular}{@{}cccc@{}}
\toprule
$w$&$w/4$&$\partial\mathcal L/\partial w^{[1]}$
 &$\partial\mathcal L/\partial w^{[4]}$\\
\midrule
$1$&$0{,}25$&$1/512=0{,}001953125$&$0{,}0625$\\
$100$&$25$&$1953{,}125$&$0{,}0625$\\
\bottomrule
\end{tabular}
\end{table}
Hai mạng đều có nơ-ron ở điểm đạo hàm kích hoạt cực đại, nhưng
građien đầu mạng rất khác nhau do trọng số. Ví dụ này chứng minh
bão hòa không phải nguyên nhân duy nhất gây khó khăn, và trọng số
lớn không tất yếu làm mọi tiền kích hoạt lớn tại một mẫu.
Cấu hình trên không bảo đảm $z^{[l]}=0$ đồng thời với mọi đầu vào.

Với $w=1$, qua đúng $r$ chuyển tiếp, tỉ số độ lớn của $\delta$
bằng $4^{-r}$. Khi $r=10$, tỉ số là
$4^{-10}=9{,}536743\times10^{-7}$. Phải đếm số chuyển tiếp thực
sự và phân biệt tỉ số này với tỉ số građien trọng số, vốn còn chứa
các đầu ra của lớp trước theo \eqref{kt:eq:parameter-gradient}.

\subsection{Vai trò của hàm mất mát tại đầu ra}

Với đầu ra $a=\sigma(z)$ và mất mát bình phương,
$\partial\mathcal L/\partial z=(a-y)a(1-a)$. Nếu bài toán phân
loại nhị phân dùng mất mát entropy chéo
\begin{equation}
 \ell(a,y)=-y\log a-(1-y)\log(1-a),\qquad y\in\{0,1\},
 \label{kt:eq:cross-entropy}
\end{equation}
thì, với $a\in(0,1)$,
\begin{equation}
 \frac{\partial\ell}{\partial z}
 =\left(-\frac ya+\frac{1-y}{1-a}\right)a(1-a)
 =-y(1-a)+(1-y)a=a-y.
 \label{kt:eq:cross-entropy-gradient}
\end{equation}
Trong đó, $y$ là nhãn và $a$ là xác suất dự đoán. Thừa số nhỏ của
đạo hàm logistic ở đầu ra đã được khử trong biểu thức giải tích;
điều này không khử các đạo hàm kích hoạt ở những lớp ẩn.
Công thức cần được tính bằng biểu thức ổn định số khi $a$ gần hai
đầu mút. Như vậy không thể đánh giá nguy cơ triệt tiêu građien chỉ
từ tên hàm kích hoạt mà bỏ qua cấu trúc hàm mất mát.

\section{Khởi tạo trọng số và lan truyền mômen bậc hai}
\label{kt:sec:initialization}

\subsection{Phá vỡ đối xứng và phạm vi của mô hình xác suất}

Nếu các nơ-ron trong cùng lớp có tham số vào giống nhau, và các
trọng số đi ra cũng đối xứng tương ứng, chúng cho cùng đầu ra và
nhận cùng cập nhật dưới một thuật toán xác định bảo toàn đối xứng.
Khởi tạo ngẫu nhiên độc lập các hàng trọng số là một cách phá vỡ
tình trạng này. Không cần làm ngẫu nhiên mọi độ lệch: đặt
$\mathbf b^{[l]}=0$ vẫn tương thích với việc phá đối xứng bằng trọng
số. Ngược lại, chỉ nói các hàng trọng số vào giống nhau chưa đủ
để kết luận chúng mãi giống nhau nếu các liên kết đi ra đã khác nhau.

Trong phân tích khởi tạo dưới đây, giả sử:
\begin{enumerate}[label=\textnormal{(K\arabic*)},leftmargin=*]
 \item Các lớp có trọng số được lấy độc lập; tại lớp $l$, các phần
 tử $W_{ij}^{[l]}$ độc lập, có trung bình bằng không và phương sai
 chung $s_l^2<\infty$.
 \item Ma trận $\mathbf W^{[l]}$ độc lập với
 $\mathbf a^{[l-1]}$. Đầu vào gốc độc lập với các phép khởi tạo.
 \item Các độ lệch bằng không và các mômen bậc hai đang xét hữu hạn.
\end{enumerate}
Các giả thiết này áp dụng tại khởi tạo. Sau cập nhật theo dữ liệu,
độc lập thống kê giữa các lớp hoặc giữa trọng số và tín hiệu vào
không còn được bảo đảm. Độ rộng các lớp được phép khác nhau.
Kỳ vọng trong mục này lấy theo phép khởi tạo và, nếu đầu vào ngẫu
nhiên, cả phân phối đầu vào; nó không phải tự động là trung bình
theo một lô dữ liệu của một mạng đã cố định.

\subsection{Một đẳng thức chính xác cho lan truyền thuận}

Điều kiện theo $\mathbf a=\mathbf a^{[l-1]}$ và dùng
$z_i^{[l]}=\sum_j W_{ij}^{[l]}a_j$, ta có
\begin{align}
 \mathbb E[z_i^{[l]}\mid\mathbf a]&=0,\notag\\
 \mathbb E[(z_i^{[l]})^2\mid\mathbf a]
 &=\sum_j a_j^2\mathbb E[(W_{ij}^{[l]})^2]
   +\sum_{j\ne k}a_ja_k\mathbb E[W_{ij}^{[l]}W_{ik}^{[l]}]\notag\\
 &=s_l^2\sum_j a_j^2.
 \label{kt:eq:conditional-second}
\end{align}
Các số hạng chéo bằng không vì trọng số khác cột độc lập và có
trung bình bằng không; ở đây các $a_j$ là những giá trị đã được
điều kiện hóa, không cần độc lập với nhau.
Lấy kỳ vọng một lần nữa cho
\begin{equation}
 \mathbb E[z_i^{[l]}]=0,\qquad
 \operatorname{Var}(z_i^{[l]})
 =s_l^2\sum_{j=1}^{n_{l-1}}\mathbb E[(a_j^{[l-1]})^2].
 \label{kt:eq:forward-moment-general}
\end{equation}
Trong đó, $\operatorname{Var}(X)=\mathbb E[X^2]-(\mathbb EX)^2$.
Nếu các đầu vào có cùng mômen bậc hai $m_{l-1}$, đặt
\begin{equation}
 m_{l-1}=\mathbb E[(a_j^{[l-1]})^2],\qquad
 q_l=\operatorname{Var}(z_i^{[l]}),\qquad
 q_l=n_{l-1}s_l^2m_{l-1}.
 \label{kt:eq:forward-moment}
\end{equation}
Ký hiệu $m$ là mômen bậc hai chưa trừ trung bình; $q$ là phương
sai tiền kích hoạt có trung bình bằng không. Nếu các $a_j$ không
có cùng mômen, vẫn dùng đẳng thức tổng
\eqref{kt:eq:forward-moment-general}.

Điểm cần giữ đúng là
\begin{equation}
 m_{l-1}=\operatorname{Var}(a_j^{[l-1]})
                +(\mathbb E[a_j^{[l-1]}])^2.
 \label{kt:eq:moment-versus-variance}
\end{equation}
Với đầu ra logistic hoặc ReLU, trung bình thường khác không; thay
$m_{l-1}$ bằng phương sai sẽ làm mất một số hạng. Nếu thêm độ lệch
ngẫu nhiên độc lập, có trung bình bằng không và phương sai
$\tau_l^2$, thì vế phải của \eqref{kt:eq:forward-moment} cộng
thêm $\tau_l^2$. Độ lệch xác định khác không còn làm dịch chuyển
trung bình và có thể phá tính đối xứng được dùng bên dưới.

\subsection{Mạng tuyến tính và nguyên lý chọn thang theo số đầu vào}

Nếu $\phi_l(z)=z$, các lớp đều rộng $n$ và $s_l^2=s^2$, thì
\eqref{kt:eq:forward-moment} cho đệ quy chính xác
\begin{equation}
 m_l=ns^2m_{l-1},\qquad m_L=(ns^2)^Lm_0.
 \label{kt:eq:linear-depth}
\end{equation}
Trong đó, $m_0$ là mômen bậc hai chung của các tọa độ đầu vào.
Nếu $0<ns^2<1$ thì $m_L\to0$; nếu $ns^2>1$ và $m_0>0$ thì
$m_L\to\infty$; nếu $ns^2=1$ thì $m_L=m_0$. Các kết luận nói về
mômen của một họ mạng ngẫu nhiên, không mô tả đầy đủ từng mạng đã
lấy mẫu. Chẳng hạn, trong trường hợp co, với $t>0$,
\[
 \mathbb P(|a_i^{[L]}|\geq t)\leq m_L/t^2\longrightarrow0,
\]
theo bất đẳng thức Markov. Đây là hệ quả xác suất của sự hội tụ
theo trung bình bình phương.

Với độ rộng thay đổi, điều kiện bảo toàn mômen từng tọa độ theo
chiều thuận của mô hình tuyến tính là
\begin{equation}
 s_l^2=\frac1{n_{l-1}}.
 \label{kt:eq:lecun}
\end{equation}
Đây là thang khởi tạo theo số đầu vào, thường gắn với khởi tạo
LeCun; việc chọn thang trọng số và chuẩn hóa dữ liệu được thảo luận
trong \cite{lecun1998}. Với $\tanh z\approx z$ gần không, công
thức có thể làm chuẩn tham chiếu ban đầu, nhưng tính phi tuyến làm
cho đẳng thức bảo toàn tuyến tính không còn áp dụng nguyên dạng.
Ngưỡng $1/n$ cũng không được chuyển nguyên sang mọi hàm kích hoạt.

\subsection{Tính đối xứng và nguồn gốc hệ số một phần hai của ReLU}

Đặt $\phi(z)=\max\{0,z\}$, gọi là hàm tuyến tính chỉnh lưu ReLU.
Ngoài (K1)--(K3), giả sử mỗi hàng trọng số có phân phối đối xứng
qua gốc. Điều kiện theo đầu vào $\mathbf a$, hai vectơ trọng số
$\mathbf W_{i,:}$ và $-\mathbf W_{i,:}$ có cùng phân phối, nên
$z_i$ và $-z_i$ cũng có cùng phân phối. Lấy trung bình theo
$\mathbf a$ vẫn giữ tính đối xứng này.

\begin{proposition}[Mômen của phần dương]
Nếu $Z$ đối xứng quanh không và $\mathbb E[Z^2]<\infty$, thì
\begin{equation}
 \mathbb E[(\max\{0,Z\})^2]=\frac12\mathbb E[Z^2].
 \label{kt:eq:relu-half}
\end{equation}
\end{proposition}
\begin{proof}
Tính đối xứng cho
$\mathbb E[Z^2\mathbf1\{Z>0\}]
=\mathbb E[Z^2\mathbf1\{Z<0\}]$.
Hai số hạng cộng lại bằng $\mathbb E[Z^2]$, vì đóng góp tại
$Z=0$ bằng không. Số hạng thứ nhất chính là vế trái cần tính.
\end{proof}
Mệnh đề vẫn đúng khi có khối lượng xác suất tại không. Trong khi
đó, đẳng thức $\mathbb P(Z>0)=1/2$ còn cần
$\mathbb P(Z=0)=0$; nói chung,
\begin{equation}
 \mathbb P(Z>0)=\mathbb P(Z<0)
       =\frac{1-\mathbb P(Z=0)}2.
 \label{kt:eq:zero-atom}
\end{equation}
Vì vậy không chứng minh \eqref{kt:eq:relu-half} bằng cách bỏ qua
giá trị bằng không của đầu ra ReLU hoặc thay tổng nhiều tích bằng
một tích đại diện.

Nếu riêng $Z\sim\mathcal N(0,q)$, $q>0$, tích phân trực tiếp cho
\begin{align}
 \mathbb E[\max\{0,Z\}]
 &=\frac1{\sqrt{2\pi q}}\int_0^\infty
               z\exp(-z^2/(2q))\,\mathrm dz
   =\sqrt{\frac q{2\pi}},\notag\\
 \operatorname{Var}(\max\{0,Z\})
 &=\frac q2-\frac q{2\pi}
   =q\left(\frac12-\frac1{2\pi}\right).
 \label{kt:eq:relu-gaussian-variance}
\end{align}
Trong đó, $\mathcal N(0,q)$ là phân phối chuẩn có phương sai $q$.
Bước tích phân dùng đạo hàm của $\exp(-z^2/(2q))$.
Như vậy hệ số $1/2$ trong \eqref{kt:eq:relu-half} thuộc mômen
bậc hai, không phải phương sai của đầu ra ReLU.

\subsection{Khởi tạo He theo chiều thuận}

Từ $m_l=q_l/2$ và \eqref{kt:eq:forward-moment}, với các lớp ReLU,
\begin{equation}
 m_l=\frac12n_{l-1}s_l^2m_{l-1}.
 \label{kt:eq:he-forward-m}
\end{equation}
Với $l\geq2$, khi lớp trước cũng dùng ReLU, tương đương
\begin{equation}
 q_l=\frac12n_{l-1}s_l^2q_{l-1},\qquad
 q_L=q_1\prod_{l=2}^{L}\left(\frac12n_{l-1}s_l^2\right).
 \label{kt:eq:he-forward-q}
\end{equation}
Để mỗi hệ số nhân bằng một, chọn
\begin{equation}
 \boxed{s_l^2=\frac2{n_{l-1}}.}
 \label{kt:eq:he-in}
\end{equation}
Trong đó, $n_{l-1}$ là số đầu vào của một nơ-ron tại lớp $l$.
Hệ số hai bù cho việc hàm chỉnh lưu giữ một nửa mômen bậc hai của
tiền kích hoạt đối xứng. Đây là cơ sở của khởi tạo He
\cite{he2015}; trong mô hình thuận vừa phát biểu, phép suy ra
không cần xấp xỉ chuẩn hoặc giới hạn bề rộng vô hạn.

Ở lớp đầu tiên với đầu vào thô có mômen $m_0$, công thức cho
$q_1=2m_0$ và $m_1=m_0$. Không có sự mâu thuẫn: đại lượng được
bảo toàn từ dữ liệu sang đầu ra ReLU là $m$, không phải một
``phương sai tiền kích hoạt lớp không'' chưa được định nghĩa.
Nếu lớp đầu ra là tuyến tính, thì $m_L=q_L$ và điều kiện thuận
$s_L^2=1/n_{L-1}$ thay cho hệ số hai khi muốn bảo toàn $m$ tại
lớp đó. Do đó cần xét riêng hàm kích hoạt của lớp đầu ra.

\paragraph{Mở rộng cho nhánh âm có hệ số khác không.}
Cho
\begin{equation}
 \phi_\alpha(z)=
 \begin{cases}z,&z\geq0,\\\alpha z,&z<0,\end{cases}
 \qquad 0\leq\alpha<1.
 \label{kt:eq:leaky}
\end{equation}
Trong đó, $\alpha$ là hệ số góc trên nửa trục âm. Với $Z$ đối
xứng, tách tích phân theo dấu như trong mệnh đề cho
\begin{equation}
 \mathbb E[\phi_\alpha(Z)^2]
   =\frac{1+\alpha^2}{2}\mathbb E[Z^2],\qquad
 s_l^2=\frac{2}{(1+\alpha^2)n_{l-1}}.
 \label{kt:eq:leaky-init}
\end{equation}
Công thức thứ hai làm hệ số nhân mômen thuận bằng một. Khi
$\alpha=0$, ta thu lại ReLU. Nếu $\alpha$ được học trong quá
trình tối ưu, lựa chọn này chỉ gắn với giá trị khởi tạo của nó.

\subsection{Lan truyền ngược: vì sao cần một xấp xỉ bổ sung}

Quan hệ đại số \eqref{kt:eq:delta-recursion} là chính xác, nhưng
việc tách kỳ vọng các thừa số trong đó không tự động đúng. Đặt
$D_j=\phi_{l-1}'(z_j^{[l-1]})$ và xét $l\geq2$. Khai triển
bình phương cho
\begin{align}
 \mathbb E[(\delta_j^{[l-1]})^2]
 &=\sum_{i=1}^{n_l}
  \mathbb E[D_j^2(W_{ij}^{[l]})^2(\delta_i^{[l]})^2]\notag\\
 &\quad+\sum_{i\ne k}
  \mathbb E[D_j^2 W_{ij}^{[l]}W_{kj}^{[l]}
                         \delta_i^{[l]}\delta_k^{[l]}].
 \label{kt:eq:backward-exact-moments}
\end{align}
Để có một mô hình thống kê đơn giản, bỏ qua các mômen chéo ở dòng
thứ hai và xấp xỉ mômen ở dòng thứ nhất bằng tích các mômen:
\begin{equation}
 \mathbb E[D_j^2(W_{ij}^{[l]})^2(\delta_i^{[l]})^2]
 \approx\mathbb E[D_j^2]\,s_l^2\,
                         \mathbb E[(\delta_i^{[l]})^2].
 \label{kt:eq:backward-factorization}
\end{equation}
Nếu các tọa độ có chung mômen, đặt
$r_l=\mathbb E[(\delta_i^{[l]})^2]$ và
$\chi_{l-1}=\mathbb E[\phi_{l-1}'(Z_{l-1})^2]$, với
$Z_{l-1}$ là tiền kích hoạt đại diện. Khi đó
\begin{equation}
 r_{l-1}\approx n_ls_l^2\chi_{l-1}r_l,
 \qquad
 r_l\approx r_L\prod_{k=l+1}^{L}
                    \left(n_ks_k^2\chi_{k-1}\right).
 \label{kt:eq:backward-moment}
\end{equation}
Ký hiệu $r_l$ là mômen bậc hai, không mặc nhiên là phương sai,
vì trung bình của građien có thể khác không. Thống kê cần dùng là
kỳ vọng của \emph{bình phương} đạo hàm, không phải bình phương
của kỳ vọng đạo hàm. Ví dụ, nếu $D$ bằng không hoặc một với xác
suất bằng nhau thì $\mathbb E[D^2]=1/2$, trong khi
$(\mathbb ED)^2=1/4$.

\paragraph{Một phản ví dụ cho giả thiết độc lập tự động.}
Xét $f(x)=Wx$, $x=1$, $\mathcal L=\tfrac12f(x)^2$ và
$W\sim\mathcal N(0,s^2)$. Đạo hàm theo đầu ra là $\delta=W$,
nên
\begin{equation}
 \mathbb E[W\delta]=\mathbb E[W^2]=s^2\ne
                    \mathbb E[W]\mathbb E[\delta]=0.
 \label{kt:eq:dependency-counterexample}
\end{equation}
Ngay tại khởi tạo, tín hiệu lỗi đã phụ thuộc trọng số thông qua
hàm mất mát. Do đó \eqref{kt:eq:backward-factorization} là một
giả thiết đóng mô hình hoặc một xấp xỉ cần đánh giá, không phải
hệ quả của việc khởi tạo độc lập các phần tử trọng số. Bản phân
tích của He và cộng sự cũng nêu các giả thiết độc lập khi suy ra
quy tắc theo chiều ngược~\cite{he2015}.

Nếu tiền kích hoạt ReLU đối xứng và không có khối lượng tại không,
thì $\chi_{l-1}=1/2$. Quy tắc giữ mômen građien từng tọa độ trong
mô hình xấp xỉ là
\begin{equation}
 \boxed{s_l^2=\frac2{n_l}.}
 \label{kt:eq:he-out}
\end{equation}
Ở đây $n_l$ là số đầu ra của lớp có ma trận $\mathbf W^{[l]}$.
Nếu quy ước đạo hàm ReLU tại không bằng không và xác suất tại
không là $p_0>0$, thì $\chi=(1-p_0)/2$, không còn đúng bằng
$1/2$. Một mạng hữu hạn có thể có tiền kích hoạt bằng không với
xác suất dương khi cả lớp trước cho vectơ không.

Khi chọn quy tắc thuận \eqref{kt:eq:he-in},
\begin{equation}
 r_{l-1}\approx\frac{n_l}{n_{l-1}}r_l,\qquad
 r_l\approx\frac{n_L}{n_l}r_L
 \label{kt:eq:width-telescope}
\end{equation}
trong chuỗi chuyển tiếp thỏa các giả thiết trên. Tỉ số độ rộng
triệt tiêu theo tích, nên bảo toàn mômen mỗi tọa độ và bảo toàn
chuẩn bình phương toàn vectơ là hai tiêu chuẩn khác nhau. Thật
vậy, $\mathbb E\|\boldsymbol\delta^{[l]}\|_2^2=n_lr_l$ nếu các
tọa độ có cùng mômen. Công thức đệ quy cho $\delta$ chỉ áp dụng
giữa các lớp có tiền kích hoạt; đạo hàm theo đầu vào gốc là
$(\mathbf W^{[1]})^T\boldsymbol\delta^{[1]}$ và không chứa đạo
hàm kích hoạt của một lớp không có thật.

\subsection{Khởi tạo Xavier và ý nghĩa của phép dung hòa}

Xét các hàm kích hoạt gần tuyến tính quanh không, có
$\phi(0)=0$ và $\phi'(0)=1$. Trong chế độ đó, lan truyền thuận
gợi ý hệ số $G_{\mathrm t}=n_{l-1}s_l^2$, còn mô hình ngược gợi
ý $G_{\mathrm n}=n_ls_l^2$. Yêu cầu giữ từng hệ số bằng một là
\begin{equation}
 s_l^2=\frac1{n_{l-1}},\qquad s_l^2=\frac1{n_l}.
 \label{kt:eq:xavier-two-conditions}
\end{equation}
Trong đó, các chỉ số $\mathrm t$ và $\mathrm n$ chỉ chiều thuận
và chiều ngược. Hai phương trình chỉ đồng thời thỏa khi
$n_{l-1}=n_l$.

Một quy tắc dung hòa là yêu cầu trung bình của hai hệ số khuếch
đại bằng một:
\begin{equation}
 \frac{G_{\mathrm t}+G_{\mathrm n}}2=1
 \quad\Longleftrightarrow\quad
 \boxed{s_l^2=\frac2{n_{l-1}+n_l}.}
 \label{kt:eq:xavier}
\end{equation}
Đây là quy tắc Xavier của Glorot và Bengio~\cite{glorot2010}.
Cách viết qua trung bình ở trên giải thích đại số của phép dung
hòa, không phải chứng minh đây là lựa chọn tối ưu duy nhất.
Nó không giữ riêng cả hai hệ số bằng một khi các độ rộng khác nhau.

Với $\tanh$, các giả thiết tuyến tính hóa có cơ sở gần không.
Tuy nhiên, $\tanh^2z\leq z^2$ cho thấy tác động phi tuyến có thể
tiếp tục làm co mômen, ngay cả khi quy tắc tuyến tính đã cân bằng.
Với logistic, $\sigma(0)=1/2$ và $\sigma'(0)=1/4$; xấp xỉ địa
phương là $\sigma(z)=1/2+z/4+o(|z|)$, chứa một số hạng hằng.
Vì vậy không áp dụng trực tiếp lập luận có trung bình bằng không
cho logistic hoặc chỉ nhân phương sai Xavier với một hệ số mà
bỏ qua số hạng hằng này.

\subsection{Chọn phân phối và kiểm tra một ví dụ số}

Một phương sai mục tiêu $s_l^2$ có thể được thực hiện bằng
\begin{equation}
 W_{ij}^{[l]}\sim\mathcal N(0,s_l^2)
 \quad\text{hoặc}\quad
 W_{ij}^{[l]}\sim\mathcal U[-a_l,a_l],\qquad a_l=\sqrt{3s_l^2}.
 \label{kt:eq:distributions}
\end{equation}
Trong đó, $\mathcal U[-a,a]$ là phân phối đều đối xứng. Công thức
cho $a_l$ được suy ra từ
\begin{equation}
 \mathbb EW=0,\qquad
 \operatorname{Var}(W)=\frac1{2a}\int_{-a}^{a}w^2\,\mathrm dw
                     =\frac{a^2}{3}.
 \label{kt:eq:uniform-variance}
\end{equation}
Do đó các dạng đều của hai quy tắc chính là
\begin{align}
 \text{Xavier:}\quad
 W_{ij}^{[l]}&\sim\mathcal U\left[
 -\sqrt{\frac6{n_{l-1}+n_l}},\sqrt{\frac6{n_{l-1}+n_l}}\right],
 \label{kt:eq:xavier-uniform}\\
 \text{He theo chiều thuận:}\quad
 W_{ij}^{[l]}&\sim\mathcal U\left[
 -\sqrt{\frac6{n_{l-1}}},\sqrt{\frac6{n_{l-1}}}\right].
 \label{kt:eq:he-uniform}
\end{align}
Phân phối chuẩn dùng \emph{độ lệch chuẩn} $s_l$, không dùng
$s_l^2$ làm độ lệch chuẩn. Hai phân phối có cùng phương sai vẫn
có thể khác nhau về các mômen bậc cao và tính chất phổ của ma trận.

Xét một lớp có $n_{l-1}=128$ và $n_l=64$. Các giá trị tính trực
tiếp từ công thức là
\begin{table}[htbp]
\centering\small
\caption{Thang khởi tạo cho ma trận có $64$ hàng và $128$ cột.}
\label{kt:tab:initialization}
\begin{tabular}{@{}lccc@{}}
\toprule
Quy tắc&$s_l^2$&$s_l$&$a_l$ của phân phối đều\\
\midrule
Theo số đầu vào, tuyến tính&$1/128$&$0{,}088388$&$0{,}153093$\\
Xavier&$1/96$&$0{,}102062$&$0{,}176777$\\
He theo chiều thuận&$1/64$&$0{,}125000$&$0{,}216506$\\
He theo chiều ngược&$1/32$&$0{,}176777$&$0{,}306186$\\
\bottomrule
\end{tabular}
\end{table}
Với Xavier trong mô hình tuyến tính,
$G_{\mathrm t}=128/96=4/3$ và $G_{\mathrm n}=64/96=2/3$.
Với He thuận trong mô hình ReLU, hệ số thuận bằng
$128/(2\cdot64)=1$, còn hệ số ngược xấp xỉ
$64/(2\cdot64)=1/2$. Với He ngược, hai hệ số lần lượt bằng $2$
và $1$. Đây là các hệ số của mômen bậc hai; hệ số của căn mômen
bậc hai là căn bậc hai của chúng. Các con số không phải kết quả
thực nghiệm và không so sánh hai quy tắc trong cùng một mô hình
phi tuyến nếu giả thiết kích hoạt khác nhau.

\subsection{Liên hệ phương sai tiền kích hoạt với xác suất bão hòa}

Nếu $Z\sim\mathcal N(0,q)$ và $c>0$, chuẩn hóa $Z/\sqrt q$ cho
\begin{equation}
 \mathbb P(|Z|\geq c)
 =2\left[1-\Phi\left(\frac c{\sqrt q}\right)\right],
 \label{kt:eq:gaussian-tail}
\end{equation}
với $\Phi$ là hàm phân phối của chuẩn tắc. Khi $q\to\infty$,
$c/\sqrt q\to0$ và $\Phi(0)=1/2$, nên xác suất này tiến tới một.
Đây là kết luận cho họ phân phối chuẩn, không phải hệ quả chỉ từ
việc phương sai tăng.

Trong khởi tạo bằng trọng số chuẩn, phát biểu chính xác là
\begin{equation}
 z_i^{[l]}\mid\mathbf a^{[l-1]}
 \sim\mathcal N\left(0,s_l^2\|\mathbf a^{[l-1]}\|_2^2\right).
 \label{kt:eq:conditional-normal}
\end{equation}
Nếu $\|\mathbf a^{[l-1]}\|_2>0$, dùng
$q=s_l^2\|\mathbf a^{[l-1]}\|_2^2$ trong
\eqref{kt:eq:gaussian-tail} để có xác suất có điều kiện. Nếu
$\mathbf a^{[l-1]}=0$, tiền kích hoạt bằng không. Sau khi bỏ điều
kiện, phân phối có thể là hỗn hợp các phân phối chuẩn có phương
sai khác nhau; nó không tự động bằng một phân phối chuẩn với
phương sai trung bình. Một xấp xỉ chuẩn không điều kiện cần thêm
lập luận, chẳng hạn sự tập trung của chuẩn bình phương đầu vào.

Để thấy phương sai lớn chưa đủ, cho $Z_n$ nhận $n$ và $-n$ với
xác suất $1/(2n)$ mỗi giá trị, nhận không với xác suất $1-1/n$.
Khi đó
\begin{equation}
 \mathbb EZ_n=0,\qquad\operatorname{Var}(Z_n)=n\to\infty,
 \qquad\mathbb P(|Z_n|>c)=\frac1n\to0\quad(n>c).
 \label{kt:eq:variance-tail-counterexample}
\end{equation}
Mômen lớn do các giá trị hiếm rất lớn, trong khi phần lớn xác suất
vẫn tập trung tại không. Ngoài ra, đầu ra của logistic và $\tanh$
bị chặn, nên không thể suy ra phương sai đầu ra của chúng bùng nổ
như trong mạng tuyến tính; sự tăng tiền kích hoạt có thể biểu hiện
bằng bão hòa thay vì đầu ra tăng vô hạn.

\subsection{Giới hạn của nguyên lý bảo toàn mômen}

Bảo toàn kỳ vọng của bình phương không bảo đảm từng tích ngẫu
nhiên có độ lớn ổn định. Một ví dụ một chiều là các trọng số độc
lập $W_k$, dấu dương hoặc âm đồng xác suất và độc lập với độ lớn,
trong đó
\[
 |W_k|=\sqrt{1/2}\quad\text{hoặc}\quad\sqrt{3/2}
 \quad\text{với xác suất }1/2\text{ mỗi trường hợp}.
\]
Đặt $P_L=\prod_{k=1}^{L}W_k$. Tính độc lập cho
\begin{equation}
 \mathbb E[P_L^2]=\prod_{k=1}^{L}\mathbb E[W_k^2]=1.
 \label{kt:eq:mean-square-product}
\end{equation}
Tuy nhiên, luật số lớn áp dụng cho $\log|W_k|$ cho
\begin{equation}
 \frac1L\log|P_L|
 \longrightarrow\mathbb E\log|W_1|
 =\frac14\log(3/4)<0
 \quad\text{gần chắc chắn}.
 \label{kt:eq:typical-product}
\end{equation}
Do đó $|P_L|\to0$ gần chắc chắn dù mômen bậc hai luôn bằng một.
Không có mâu thuẫn: một số biến cố hiếm giữ đóng góp lớn vào kỳ
vọng, và không thể tự ý đổi giới hạn với kỳ vọng.

Ở mức ma trận, bảo toàn một đại lượng trung bình cũng không giữ
mọi hướng. Với $0<\epsilon<1$, lấy
\begin{equation}
 \mathbf J=\operatorname{diag}(\sqrt{2-\epsilon^2},\epsilon).
 \quad\text{Khi đó}\quad
 \frac12\operatorname{tr}(\mathbf J^T\mathbf J)=1,
 \qquad\kappa_2(\mathbf J)=\frac{\sqrt{2-\epsilon^2}}{\epsilon}.
 \label{kt:eq:singular-counterexample}
\end{equation}
Trong đó, $\operatorname{tr}$ là tổng các phần tử đường chéo và
$\kappa_2$ là tỉ số giá trị kỳ dị lớn nhất với nhỏ nhất. Khi
$\epsilon\to0$, số điều kiện tăng vô hạn dù trung bình bình
phương các giá trị kỳ dị bằng một. Vì vậy khởi tạo theo mômen là
một tiêu chuẩn kiểm soát thang đo ban đầu, không phải một chứng
minh mọi hướng của phép lan truyền đều ổn định.

\section{Các biện pháp hỗ trợ và lưu ý đối với mạng giải phương trình vi phân}
\label{kt:sec:practical}

\subsection{Chuẩn hóa đầu vào và tính nhất quán của đạo hàm}

Với từng tọa độ, có thể dùng biến đã chuẩn hóa
\begin{equation}
 \widehat x_j=\frac{x_j-\mu_j}{s_j},\qquad s_j>0.
 \label{kt:eq:input-scaling}
\end{equation}
Ở đây $\mu_j$ là mốc dịch chuyển và $s_j$ là thang đo, được chọn
từ miền bài toán hoặc tập huấn luyện rồi giữ nhất quán khi đánh
giá. Phép biến đổi giúp tránh sự chênh lệch đơn vị quá lớn giữa
các tọa độ; nó không tự động bảo đảm tiền kích hoạt của mọi lớp
nằm ngoài vùng bão hòa.

Trong mạng tích hợp thông tin vật lý, nếu
$u(\mathbf x)=\widehat u(\widehat{\mathbf x})$, quy tắc dây chuyền
đòi hỏi
\begin{equation}
 \partial_{x_j}u=s_j^{-1}\partial_{\widehat x_j}\widehat u,
 \qquad
 \partial_{x_jx_j}u=s_j^{-2}
                   \partial_{\widehat x_j\widehat x_j}\widehat u.
 \label{kt:eq:pde-scaling}
\end{equation}
Các hệ số xuất hiện do
$\partial\widehat x_j/\partial x_j=1/s_j$ và lấy đạo hàm lần
nữa. Bỏ chúng khi lập phần dư sẽ làm thay đổi phương trình đang
giải. Nếu cả đầu ra cũng được đổi thang thì phải đưa thêm hệ số
tương ứng vào các đạo hàm.

\subsection{Chuẩn hóa theo lô}

Với một tọa độ tiền kích hoạt trên lô gồm $B$ mẫu, phép chuẩn hóa
theo lô của Ioffe và Szegedy~\cite{ioffe2015} có dạng
\begin{align}
 \mu_B&=\frac1B\sum_{p=1}^{B}z_p,&
 v_B&=\frac1B\sum_{p=1}^{B}(z_p-\mu_B)^2,\notag\\
 \widehat z_p&=\frac{z_p-\mu_B}{\sqrt{v_B+\epsilon_B}},&
 y_p&=\gamma\widehat z_p+\beta.
 \label{kt:eq:batch-normalization}
\end{align}
Trong đó, $\epsilon_B>0$ tránh mẫu số bằng không; $\gamma$ và
$\beta$ là các tham số được học; $\mu_B,v_B$ là thống kê của tọa
độ trên lô hiện tại. Theo định nghĩa,
\begin{equation}
 \frac1B\sum_p\widehat z_p=0,\qquad
 \frac1B\sum_p\widehat z_p^2=\frac{v_B}{v_B+\epsilon_B}\leq1.
 \label{kt:eq:bn-moments}
\end{equation}
Vì $\epsilon_B>0$, phương sai thực nghiệm sau chuẩn hóa không
nhất thiết đúng bằng một; sau phép biến đổi với $\gamma,\beta$,
tín hiệu vẫn có thể đi vào vùng bão hòa.

Lời giải thích ban đầu dựa trên giảm biến động phân phối tín hiệu
giữa các lớp không nên được trình bày như cơ chế duy nhất đã được
chứng minh. Santurkar và cộng sự~\cite{santurkar2018} phân tích
thêm tác động của phép chuẩn hóa lên tính trơn của bài toán tối
ưu. Trong mạng tính phần dư phương trình vi phân, còn cần chú ý
$\mu_B,v_B$ phụ thuộc cả lô: đầu ra tại một điểm khi huấn luyện
có thể phụ thuộc các điểm khác trong lô. Khi lấy đạo hàm theo tọa
độ, phải xác định rõ dùng thống kê cố định hay thống kê phụ thuộc
lô, thay vì mặc nhiên xem đây là một ánh xạ từng điểm không đổi.

\subsection{Giới hạn chuẩn građien}

Cho $\mathbf g=\nabla_{\boldsymbol\theta}\mathcal L$ và ngưỡng
$C>0$, đặt
\begin{equation}
 \widetilde{\mathbf g}=
 \begin{cases}
  \mathbf g,&\|\mathbf g\|_2\leq C,\\
  C\mathbf g/\|\mathbf g\|_2,&\|\mathbf g\|_2>C.
 \end{cases}
 \label{kt:eq:clipping}
\end{equation}
Trong đó, $\boldsymbol\theta$ ghép các tham số được tối ưu.
Phép biến đổi giữ hướng khi $\mathbf g\ne0$ và bảo đảm
$\|\widetilde{\mathbf g}\|_2\leq C$. Với hạ dốc đơn thuần,
\begin{equation}
 \boldsymbol\theta_{k+1}
 =\boldsymbol\theta_k-\eta_k\widetilde{\mathbf g}_k
 \quad\Longrightarrow\quad
 \|\boldsymbol\theta_{k+1}-\boldsymbol\theta_k\|_2\leq\eta_k C.
 \label{kt:eq:clipped-step}
\end{equation}
Ở đây $k$ là bước tối ưu và $\eta_k>0$ là chiều dài bước. Cận này
theo trực tiếp tính thuần nhất của chuẩn; không chuyển nguyên
sang thuật toán có bộ nhớ hoặc tiền điều kiện mà không xét công
thức cập nhật của thuật toán đó. Giới hạn chuẩn được dùng để
kiểm soát građien lớn, tiêu biểu trong nghiên cứu của Pascanu và
cộng sự~\cite{pascanu2013}; nó không phục hồi thông tin đã mất do
građien triệt tiêu và không tự bảo đảm hội tụ.

\subsection{Kết nối tắt và lựa chọn hàm kích hoạt}

Trong một khối dư có các trạng thái cùng chiều,
\begin{equation}
 \mathbf a^{[l+1]}=\mathbf a^{[l]}+F_l(\mathbf a^{[l]}),\qquad
 \frac{\partial\mathbf a^{[l+1]}}{\partial\mathbf a^{[l]}}
 =\mathbf I+DF_l(\mathbf a^{[l]}).
 \label{kt:eq:residual-block}
\end{equation}
Trong đó, $F_l$ là ánh xạ khả vi và $DF_l$ là ma trận Jacobi của
nó. Đây là cơ chế kết nối tắt trong học phần dư~\cite{he2016}.
Nếu $\|DF_l\|_2\leq\rho<1$, bất đẳng thức tam giác thuận và
ngược cho mọi $\mathbf v$:
\begin{equation}
 (1-\rho)\|\mathbf v\|_2
 \leq\|(\mathbf I+DF_l)\mathbf v\|_2
 \leq(1+\rho)\|\mathbf v\|_2.
 \label{kt:eq:residual-bound}
\end{equation}
Công thức giải thích lợi ích của một khối có phần hiệu chỉnh nhỏ.
Tuy nhiên, qua nhiều khối, tích các cận vẫn có thể co hoặc tăng;
kết nối tắt không phải bảo đảm tuyệt đối loại bỏ mọi bệnh lý
građien. Đặc biệt, nó khác với việc ngẫu nhiên loại bỏ nơ-ron hoặc
kết nối trong quá trình huấn luyện.

ReLU không bão hòa trên nửa trục dương, nhưng đạo hàm bằng không
trên nửa trục âm. Với một mạng truyền thẳng dùng ReLU và đầu ra
afin, hàm theo đầu vào là afin liên tục từng khúc; đạo hàm bậc
hai bằng không bên trong mỗi miền tuyến tính. Vì vậy việc chọn
ReLU thuận lợi cho lan truyền đạo hàm bậc nhất không tự nó làm
ReLU phù hợp cho phần dư mạnh của phương trình bậc hai.
Hàm $\tanh$ đủ trơn cho các đạo hàm cổ điển ở mọi bậc hữu hạn,
nhưng vẫn có thể bão hòa. Lựa chọn kích hoạt phải xét đồng thời
độ trơn mà phương trình đòi hỏi, thang đầu vào, kiến trúc và mô
hình khởi tạo.

Trong đánh giá số, nên báo cáo phân phối tiền kích hoạt, tỉ lệ
bão hòa theo ngưỡng đã công bố, độ lớn građien từng lớp và độ
lớn cập nhật thực tế. Các đại lượng này cần được theo dõi cả lúc
khởi tạo lẫn trong huấn luyện. Chúng giúp kiểm tra các giả thiết
của phân tích thay vì coi một tên phương pháp khởi tạo là bằng
chứng rằng toàn bộ quá trình tối ưu đã ổn định.
